
\begin{register}{H}{tselect}{0x7A0}\label{regdesc:TSELECT}
\regfield{(reserved)}{29}{3}{{0x00000000}}%
\regfield{tselect}{3}{0}{{0x0}}%
\reglabel{Reset}\regnewline%
\begin{regdesc}\begin{reglist}
\label{fielddesc:TSELECT}\item [tselect] 触发器单元编号 (0-7)。（读/写）
\end{reglist}\end{regdesc}
\end{register}

\begin{register}{H}{tdata1}{0x7A1}\label{regdesc:TDATA1}
\regfield{type}{4}{28}{{0x2}}%
\regfield{dmode}{1}{27}{{0}}%
\regfield{data}{27}{0}{{0x3e00000}}%
\reglabel{Reset}\regnewline%
\begin{regdesc}\begin{reglist}
\label{fielddesc:TDATA1TYPE}\item [type] 触发器类型。（只读） \newline 仅支持匹配类型 (0x2)，此域保留。
\label{fielddesc:TDATA1DMODE}\item [dmode] 如果某触发器正在被调试器使用，则此域置为 1.（读/写 *）
\begin{itemize}
    \setlength\itemsep{-1em}
    \item 0：在调试模式和机器模式下都能写入 tdata1 和 tdata2
    \item 1：只有在调试模式下才能写入 tdata1 和 tdata2。其他模式下的写操作将被忽略。
\end{itemize}
\textit{*说明：仅支持调试模式下的写操作。}
\label{fielddesc:TDATA1DATA}\item [data] 保存抽象 tdata1 的内容。（读/写） \newline 由于仅支持匹配类型 (0x2) 触发器，此域将始终被解读为 \hyperref[regdesc:MCONTROL]{mcontrol} 的域。
\end{reglist}\end{regdesc}
\end{register}

\begin{register}{H}{tdata2}{0x7A2}\label{regdesc:TDATA2}
\regfield{tdata2}{32}{0}{{0x00000000}}%
\reglabel{Reset}\regnewline%
\begin{regdesc}\begin{reglist}
\label{fielddesc:TDATA2}\item [tdata2] 保存抽象 tdata2 的内容。（读/写） \newline 由于仅支持匹配类型 (0x2) 触发器，此域将始终被解读为 \hyperref[regdesc:MADDRESS]{maddress}。
\end{reglist}\end{regdesc}
\end{register}

\begin{register}{H}{tcontrol}{0x7A5}\label{regdesc:TCONTROL}
\regfield{(reserved)}{24}{8}{{0x000000}}%
\regfield{mpte}{1}{7}{{0}}%
\regfield{(reserved)}{3}{4}{{0}}%
\regfield{mte}{1}{3}{{0}}%
\regfield{(reserved)}{3}{0}{{0}}%
\reglabel{Reset}\regnewline%
\begin{regdesc}\begin{reglist}
\label{fielddesc:TCONTROLMPTE}\item [mpte] 机器模式下前一个触发器使能域。（读/写）
\begin{itemize}
    \setlength\itemsep{-1em}
    \item 当 CPU 在机器模式下进入异常，\hyperref[fielddesc:TCONTROLMTE]{mte} 的值会自动写入此域。
    \item 当 CPU 执行 MRET，此域的值会返回 \hyperref[fielddesc:TCONTROLMTE]{mte}，此域变为 0。
\end{itemize}
\label{fielddesc:TCONTROLMTE}\item [mte] 机器模式下触发器使能域。（读/写）
\begin{itemize}
    \setlength\itemsep{-1em}
    \item 当 CPU 在机器模式下进入异常，此域的值会自动写入 \hyperref[fielddesc:TCONTROLMPTE]{mpte}，然后此域变为 0，并且 \hyperref[fielddesc:MCONTROLACTION]{action}=0 的触发器被全局禁用。
    \item 当 CPU 执行 MRET，\hyperref[fielddesc:TCONTROLMPTE]{mpte} 的值会自动返回此域。
\end{itemize}
\end{reglist}\end{regdesc}
\end{register}

\begin{register}{H}{mcontrol}{0x7A1}\label{regdesc:MCONTROL}
\regfield{(reserved)}{4}{28}{{0x2}}%
\regfield{dmode}{1}{27}{{0}}%
\regfield{(reserved)}{6}{21}{{0x1f}}%
\regfield{hit}{1}{20}{{0}}%
\regfield{(reserved)}{4}{16}{{0}}%
\regfield{action}{4}{12}{{0}}%
\regfield{(reserved)}{1}{11}{{0}}%
\regfield{match}{4}{7}{{0}}%
\regfield{m}{1}{6}{{0}}%
\regfield{(reserved)}{2}{4}{{0}}%
\regfield{u}{1}{3}{{0}}%
\regfield{execute}{1}{2}{{0}}%
\regfield{store}{1}{1}{{0}}%
\regfield{load}{1}{0}{{0}}%
\reglabel{Reset}\regnewline%
\begin{regdesc}\begin{reglist}
\label{fielddesc:MCONTROLDMODE}\item [dmode] 与 \hyperref[regdesc:TDATA1]{tdata1} 的 \hyperref[fielddesc:TDATA1DMODE]{dmode} 一致。
\label{fielddesc:MCONTROLHIT}\item [hit] 如果选定的触发器之前触发过，则此域为 1。（读/写） \newline 此域必须手动清零。
\label{fielddesc:MCONTROLACTION}\item [action] 配置选定的触发器在触发时进行以下操作。（读/写）
\newline 有效选项为：
\begin{itemize}
    \setlength\itemsep{-1em}
    \item 0x0：引起断点异常
    \item 0x1：进入调试模式（仅当 dmode = 1 时有效）
\end{itemize}
\textit{说明：写入无效数值会导致此域变为默认值 0x0。}
\label{fielddesc:MCONTROLMATCH}\item [match] 配置触发器进行数据/指令地址的下匹配操作。（读/写）
\newline 有效选项为：
\begin{itemize}
    \setlength\itemsep{-1em}
    \item 0x0：严格字节匹配，即与访问中某个字节对应的地址必须严格匹配 \hyperref[regdesc:MADDRESS]{maddress} 的值。
    \item 0x1：NAPOT 匹配，即访问中至少有一个字节处于  \hyperref[regdesc:MADDRESS]{maddress} 中规定的 NAPOT 区域。
\end{itemize}
\textit{说明：写入超过最大值的数值会被裁剪为最大值 0x1。}
\label{fielddesc:MCONTROLM}\item [m] 置位使选定的触发器在机器模式下操作。（读/写）
\label{fielddesc:MCONTROLU}\item [u] 置位使选定的触发器在用户模式下操作。（读/写）
\label{fielddesc:MCONTROLEXECUTE}\item [execute] 置位使选定的触发器在 CPU 执行具有匹配的虚拟地址的指令之前触发。（读/写）
\label{fielddesc:MCONTROLSTORE}\item [store] 置位使选定的触发器在 CPU 执行具有匹配的数据地址的存储器写操作之前触发。（读/写）
\label{fielddesc:MCONTROLLOAD}\item [load] 置位使选定的触发器在 CPU 执行具有匹配的数据地址的存储器读操作之前触发。（读/写）
\end{reglist}\end{regdesc}
\end{register}

\begin{register}{H}{maddress}{0x7A2}\label{regdesc:MADDRESS}
\regfield{maddress}{32}{0}{{0x00000000}}%
\reglabel{Reset}\regnewline%
\begin{regdesc}\begin{reglist}
\label{fielddesc:MADDRESS}\item [maddress] 选定的触发器执行匹配操作时使用的地址。（读/写） \newline 当 \hyperref[regdesc:MCONTROL]{mcontrol} 中的 \hyperref[fielddesc:MCONTROLMATCH]{match}=1 时由 NAPOT 解码。
\end{reglist}\end{regdesc}
\end{register}
